% =====================================================================
% CHAPTER 8: Thesis Management and Finance (new department template,
% RESTRUCTURE_PLAN_V2.md)
% Assembled from:
%   - old mainmatter/ch6.tex (Introduction, Management of the Thesis,
%     Time Frame table + prose, Overall Budget table + prose,
%     Component/Software Budget prose)
%   - old mainmatter/ch1.tex 1.5 Project Planning + 1.5.1 Work Plan
%     (Gantt figure fig:gantt, moved here per the ownership map, never
%     duplicated in the new Chapter 1)
%   - old appendices/app_c.tex "Authorship Contribution Statement"
%     (CRediT content, recast as a table, tab:credit) and "Communication"
%     (folded into 8.2 as supervision cadence)
%   - new: 8.4 Resource Planning, written from facts already established
%     in old ch6, old app_a "Software and Hardware Environment", and the
%     cross-references already fixed by ch3/ch4/ch5/ch7 in mainmatter2/.
%     No number in this file is invented; every figure traces to one of
%     the sources above.
% POs demonstrated: PO9 (individual and teamwork), PO11 (project
% management and finance).
% =====================================================================
\chapter{Thesis Management and Finance}
\label{ch:management}

\section{Introduction}
This chapter records how we managed the thesis: the schedule we followed,
the resources we drew on, what we spent, and how we split the work
between the two of us. Paper collection and the initial literature
review began during the fourth-year odd semester (the seventh semester),
and the work ran from July 2025 to July 2026 across both semesters. Model training and the experimental studies
of Chapter~\ref{ch:experimental} used rented cloud GPUs, and we also
bought a stereo camera and a Jetson Orin Nano board for real-data
capture and on-device deployment testing. Section~\ref{sec:8-2}
describes how we managed the work across eleven overlapping phases and
how we coordinated with our supervisor; Section~\ref{sec:8-3} reports
the overall project expenditure; Section~\ref{sec:8-4} sets out how we
planned compute, data, software, and our own time; and
Section~\ref{sec:8-5} states each author's individual contribution.

\FloatBarrier
\section{Management of the Thesis}
\label{sec:8-2}
We divided the work into phases with defined outputs and let some phases
overlap deliberately so that activities could run in parallel. The
literature review guided our choice of candidate architectures, and the
setup and preliminary training that followed established the
experimental pipeline. We kept capturing real-camera data while
refining the architecture through the ablation campaign of
Chapter~\ref{ch:design}, then carried the selected design through the
benchmarking, three-dimensional reconstruction, and final evaluation of
Chapter~\ref{ch:results} before consolidating the results in this book.
This overlap let data preparation, experimentation, and writing proceed
together, since design decisions still depended on experimental
evidence that only became available as the ablation campaign
progressed.

We met with our supervisor regularly throughout the project to review
progress, discuss design decisions, and plan the next steps. Milestone
reviews lined up with the end of the literature study, the architecture
design, the ablation campaign, and the production training run, and the
supervisor's feedback at each of those points shaped the direction of
the work that followed. We also prepared the final manuscript and the
defense presentation in consultation with the supervisor.

\subsection{Time Frame / Gantt Chart}
\label{sec:8-2-1}
The full project schedule ran from July 2025 to July 2026, split across
the fourth-year odd semester and the fourth-year even semester. The odd
semester covered topic selection, the literature review, the thesis
proposal, and dataset collection and processing; these four phases
fixed the topic, the research gap, and the datasets that the even
semester's experimental work then relied on. The even semester covered
the remaining seven phases, from model selection and setup through
final thesis writing. Table~\ref{tab:timeframe} lists the periods and
principal outputs of all eleven phases, and Figure~\ref{fig:gantt}
shows the corresponding Gantt chart of their overlap.

\begin{table}[!htbp]
\centering
\caption{Eleven project phases from July 2025 to July 2026.}
\label{tab:timeframe}
\small
\begin{tabular}{p{4.2cm}p{3.0cm}p{6.3cm}}
\toprule
Phase & Period & Output\\
\midrule
1. Topic selection & Jul. 2025 & Chosen thesis topic and scope\\
2. Literature review & Aug. to Sep. 2025 & Surveyed paper corpus and research gap\\
3. Thesis proposal & Sep. to Oct. 2025 & Approved proposal\\
4. Dataset collection and processing & Oct. to Nov. 2025 & Prepared training and evaluation datasets\\
5. Model selection and setup & Dec. 2025 to Jan. 2026 & Selected baseline models and a working experimental environment\\
6. Pre-training and fine-tuning & Dec. 2025 to Mar. 2026 & Initial checkpoints and a stable training procedure\\
7. Real-data capture & Jan. to Mar. 2026 & Calibrated and rectified stereo image pairs\\
8. Ablation studies & Feb. to Apr. 2026 & Evidence for the final architectural and training choices\\
9. Benchmarking and 3D reconstruction & Mar. to May 2026 & Accuracy, efficiency, and reconstruction results\\
10. Result analysis and evaluation & Apr. to Jun. 2026 & Consolidated in-domain, zero-shot, robustness, and deployment findings\\
11. Thesis writing and documentation & May to Jul. 2026 & Final thesis book and supporting documentation\\
\bottomrule
\end{tabular}
\end{table}

\begin{figure}[H]
\centering
\includegraphics[width=\textwidth]{fig_b_1_gantt_fixed}
\caption{Gantt chart of the eleven project phases from July 2025 to July 2026.}
\label{fig:gantt}
\label{fig:gantt6}
\end{figure}

\FloatBarrier
\section{Overall Budget}
\label{sec:8-3}
Table~\ref{tab:expenses} lists the approximate project expenditure.
Cloud computation made up about one third of the total: the A100
training workload cost approximately USD 100, and the other GPU
instances used for ablation and evaluation cost approximately USD 40.
The stereo camera and Jetson Orin Nano together cost about BDT 35,000,
which is approximately USD 280 at the conversion rate used in this
record.

\begin{table}[!htbp]
\centering
\caption{Project expenditure in USD and BDT across hardware and cloud computing.}
\label{tab:expenses}
\small
\begin{tabular}{p{8.4cm}rr}
\toprule
Item & Cost (USD) & Cost (BDT)\\
\midrule
Cloud A100 usage for training & 100 & 12,500\\
Other cloud GPUs for ablation and evaluation & 40 & 5,000\\
Stereo camera and Jetson Orin Nano & 280 & 35,000\\
Software licences & 0 & 0\\
\midrule
Total & 420 & 52,500\\
\bottomrule
\end{tabular}
\end{table}

\subsection{Component and Software Budget}
The budget in Table~\ref{tab:expenses} splits into hardware, cloud
computing, and software.

Hardware covered the stereo camera and the Jetson Orin Nano, used for
real-data acquisition and on-device testing, at a combined cost of
about BDT 35,000 (USD 280). Both count as development hardware for
validating the pipeline, not as the deployed system itself.

Cloud computing cost about USD 140 in total: roughly USD 100 on A100
instances for the main training run, and about USD 40 on other GPU
instances for ablation studies, intermediate checks, and final
evaluation. These are usage totals accumulated over the whole
development period, not the rate for a single day, and they will not
match a provider's current hourly tariff, since prices vary by
accelerator, region, and storage allocation.

The software stack, described in Chapter~\ref{ch:implementation}, uses
open-source tools and cost nothing in licence fees. We limited cloud
use to development; inference runs locally on the Jetson device, so
deployment adds no recurring per-frame cloud charge, only ordinary
hardware upkeep and power.

\FloatBarrier
\section{Resource Planning}
\label{sec:8-4}
Beyond the money recorded in Section~\ref{sec:8-3}, the thesis depended
on planning compute, data, software, and our own time so that each
phase of Table~\ref{tab:timeframe} had what it needed when it needed
it.

Compute split across three tiers by purpose. Local development and the
latency benchmark ran on a single RTX 3050 laptop GPU with roughly 4
gigabytes of memory, which gave us an always-available environment for
writing and debugging the training and evaluation code before
committing anything to a paid run. The production training run of
Chapter~\ref{ch:implementation} used a single rented cloud A100 with 80
gigabytes of memory, sized so the network and its training batch fit
comfortably for the one full run recorded in Table~\ref{tab:expenses}.
Ablation studies, intermediate checks, and evaluation instead ran on
cheaper rented cloud T4 instances, which cost about a tenth of the A100
workload (Table~\ref{tab:expenses}) and let us test many architectural
variants in Chapter~\ref{ch:design} without paying A100 rates for runs
that did not need the full card. We committed a design to the A100 tier
only after the T4 ablations had already validated it, so the expensive
compute trained a design we had confirmed cheaply, rather than
searching for that design at A100 prices. The Jetson Orin Nano board,
bought alongside the stereo camera in Table~\ref{tab:expenses}, served
a different role again: it is the deployment target itself, and
Chapter~\ref{ch:results} reports the on-device latency measured
directly on it rather than projected from a desktop GPU.

Data planning followed the same tiering. Chapter~\ref{ch:experimental}
describes the SceneFlow synthetic corpus and the Middlebury 2014
benchmark that trained and evaluated the network, alongside the
real-camera pairs captured with the calibrated binocular rig also
described there. The synthetic and benchmark data came from public
releases, so the data-collection cost we carried directly was the
camera hardware itself, already counted in Table~\ref{tab:expenses},
together with the time spent calibrating and rectifying our own
captures.

Software cost nothing beyond the hardware above. We built the training
and evaluation code on PyTorch 2.11 with CUDA 12.8 on Python 3.10
(Chapter~\ref{ch:implementation}), and wrote this book in LaTeX; both
toolchains are open source, which kept the software-licence row of
Table~\ref{tab:expenses} at zero throughout the project.

Finally, we planned our own time against the phase table of
Section~\ref{sec:8-2-1}. Pre-training, fine-tuning, and real-data
capture ran together over the first quarter of the schedule; ablation
studies and benchmarking followed as checkpoints matured; and we
reserved the final two months mostly for writing, so that documenting
the thesis did not compete with active experimentation for our
attention. Section~\ref{sec:8-5} records how this time was split
between the two of us.

\FloatBarrier
\section{Individual and Team Work Contribution}
\label{sec:8-5}
The thesis was carried out by a two-person team, Md. Raihanul Haque
Rahi (Roll 2008011) and Nazib Abrar (Roll 2008026), under the
supervision of Md Zunaid Hossen (Lecturer, Department of Mechatronics
Engineering), who guided the work as described in
Section~\ref{sec:8-2} but is not a co-author of this thesis. Rahi
led the architecture and training side of the work: proposing the
tile-based design, writing the model and training code, and running
the investigation and validation that the ablation campaign of
Chapter~\ref{ch:design} depended on. Abrar led the evaluation and
tooling side: building the evaluation pipeline, curating the real-camera
and benchmark data described in Chapter~\ref{ch:experimental}, and
running the validation studies that fed Chapter~\ref{ch:results}.

Both authors worked across most categories rather than staying inside a
single role, which is typical of a two-person thesis team where every
phase in Table~\ref{tab:timeframe} needs both people at once. Table
\ref{tab:credit} records each author's primary responsibilities
against the CRediT (Contributor Roles Taxonomy) categories.

\begin{table}[!htbp]
\centering
\caption{Primary CRediT responsibilities of each author.}
\label{tab:credit}
\small
\begin{tabular}{p{4.0cm}p{9.5cm}}
\toprule
Author & Primary CRediT responsibilities\\
\midrule
Md. Raihanul Haque Rahi (Roll 2008011) & Conceptualization, Methodology, Software (architecture and training code), Investigation, Validation, Writing (original draft), Visualization\\
Nazib Abrar (Roll 2008026) & Conceptualization, Methodology, Software (evaluation and tooling), Data curation, Investigation, Validation, Writing (review and editing), Visualization\\
\bottomrule
\end{tabular}
\end{table}

Both authors reviewed and approved the final manuscript.

\FloatBarrier
\section{Summary}
\label{sec:8-6}
This chapter recorded how we managed and financed the thesis. The
overall schedule ran from July 2025 to July 2026 and consisted of
eleven overlapping phases, from topic selection and the literature
review through model setup, training, and final documentation,
punctuated by milestone reviews with our supervisor. Total expenditure was about
USD 420 (BDT 52,500): USD 140 for cloud GPU usage and USD 280 for the
stereo camera and Jetson Orin Nano. We planned compute across three
tiers (an RTX 3050 laptop GPU for local development, cloud T4 instances
for ablations, and a single cloud A100 for production training), used
open-source software throughout at no licence cost, and split the work
so that architecture and training fell primarily to Rahi and evaluation
and tooling fell primarily to Abrar, with both authors reviewing every
stage of the manuscript. Because inference runs locally on the Jetson
device, deployment itself adds no recurring cloud charge beyond this
recorded development cost.

% SOURCES: mainmatter/ch6.tex (all sections: 6.1 Introduction, 6.2
%   Management of the Thesis + 6.2.1 Time Frame with tab:timeframe,
%   6.3 Overall Budget with tab:expenses, 6.4 Component/Software Budget,
%   6.5 Summary); mainmatter/ch1.tex 1.5 Project Planning + 1.5.1 Work
%   Plan (Gantt figure fig:gantt / fig_b_1_gantt_fixed, both old labels
%   fig:gantt and fig:gantt6 preserved on the single consolidated
%   figure); appendices/app_c.tex "Authorship Contribution Statement"
%   (recast as tab:credit) and "Communication" (folded into 8.2).
%   Resource Planning (8.4) numbers sourced from appendices/app_a.tex
%   "Software and Hardware Environment" (RTX 3050 ~4 GB, cloud A100 80
%   GB, cloud T4, Jetson Orin Nano, PyTorch 2.11/CUDA 12.8/Python 3.10)
%   plus the cost ratios already in tab:expenses; no new numeric claim
%   introduced.
% DROPPED: none. All content from the four source files above is
%   represented, either verbatim (tables, figure) or humanized (prose).
% MOVED-AWAY: none owned by another chapter; dataset details, rig
%   specifications, and the Jetson Orin Nano hardware description stay
%   with ch:experimental (owner ch5) and are only cross-referenced here.
